\documentclass[11pt,a4paper]{article}
% Shared preamble for RuPhO English problem statements (match T1 2026 style).
% Use: \input{latex/rupho_en_preamble.tex} after \documentclass.
\usepackage[margin=1in]{geometry}
\usepackage{amsmath,amssymb}
\usepackage{graphicx}
\usepackage{float}
\usepackage{enumitem}
\usepackage{microtype}
\usepackage{caption}
\usepackage{tabularx}
\usepackage{hyperref}

\captionsetup{font=small,labelfont=bf,skip=6pt}

\setlist[itemize]{leftmargin=1.2cm}
\setlist[enumerate]{leftmargin=1.2cm}

% One outer box per subproblem: [id | statement | points] (no inner rules)
\newcommand{\problembox}[3]{%
  \par\addvspace{0.42em}%
  \noindent
  \setlength{\fboxsep}{7.5pt}%
  \setlength{\fboxrule}{0.95pt}%
  \fbox{%
    \begin{minipage}{\dimexpr\linewidth-2\fboxsep-2\fboxrule\relax}
    \begin{tabularx}{\linewidth}{@{}>{\raggedright\arraybackslash\bfseries}p{2.1em} @{\hspace{0.9em}} >{\raggedright\arraybackslash}X @{\hspace{0.9em}} >{\raggedleft\arraybackslash\bfseries}p{2.4em}@{}}
    #1 & #2 & #3
    \end{tabularx}
    \end{minipage}%
  }%
  \par\addvspace{0.32em}%
}


\title{\textbf{Mirage}}
\author{\normalsize X17 (2017) --- English translation}
\date{}
\begin{document}
\maketitle

On a hot summer day, when the air warmed up to $t_1=35~^{\circ}\mathrm C$, the temperature of the asphalt in the sun and the thin layer of air above it was $t_2=65~^{\circ}\mathrm C$. When walking along the highway in such dry weather, “puddles” are often visible on the road. The refractive index of air depends linearly on air density: $n=1+k\rho$. At $t_0=0~^{\circ}\mathrm C$, $n_0=1{.}000292$. Consider that a person's height is $H=1{.}8~\text{m}$.

\begin{figure}[H]
    \centering
    \includegraphics[width=0.72\textwidth]{T6_2_EN_assets/fig1.jpg}
    \caption{Figure 1.}
\end{figure}

\problembox{A1}{At what minimum distance $L_{min}$ can a “puddle” be seen? Express the answer in terms of $H$, $n_0$, $T_0$, $T_1$, $T_2$ and calculate the resulting value.}{5.00}

\vspace{1em}
\noindent\textit{Source:} \url{https://pho.rs/p/3942}

\end{document}